% Exam questions with solutions — Lattice Field Theory
% Location: exam_questions_wsoln.tex
% Compile: from LFT_lectures/ run `latexmk -pdf exam_questions_wsoln.tex` or `xelatex exam_questions_wsoln.tex`

\documentclass[11pt]{article}

% Encoding and fonts
\usepackage[utf8]{inputenc}
\usepackage[T1]{fontenc}
\usepackage{lmodern}

% Math
\usepackage{amsmath,amssymb,amsthm}

% Graphics and colors
\usepackage{graphicx}
\usepackage{xcolor}
\usepackage{caption}
\usepackage{placeins}

% TikZ for diagrams
\usepackage{tikz}
\usetikzlibrary{calc, decorations.pathreplacing}

% Better typography
\usepackage{microtype}

% Page layout
\usepackage[left=1in, right=1in, top=1in, bottom=1in]{geometry}

% Hyperlinks
\usepackage[hidelinks]{hyperref}

% Lists
\usepackage{enumitem}

% Header: page style
\usepackage{fancyhdr}
\pagestyle{fancy}
\fancyhf{} % clear header and footer
\lhead{\textit{LFT : Exam Questions (with Solutions)}}
\rhead{\textit{S. Singh}}
\renewcommand{\headrulewidth}{0.4pt}
\fancyfoot[C]{\thepage} % Add centered page number in footer

% Difficulty tag
\newcommand{\diff}[1]{\textcolor{blue!60!black}{\footnotesize\textsf{[#1]}}\quad}

% Question and solution environments
\newtheorem{question}{Question}
\newenvironment{solution}{\par\smallskip\noindent\textit{Solution.}\ }{\hfill$\blacksquare$\par\smallskip}

% Metadata
\title{Lattice Field Theory — Exam Questions with Solutions}
\author{Simran Singh}
\date{\today}

\begin{document}

\maketitle

\noindent Each chapter below has three questions of increasing difficulty:
\textcolor{blue!60!black}{\textsf{[Easy]}}, \textcolor{blue!60!black}{\textsf{[Intermediate]}}, and \textcolor{blue!60!black}{\textsf{[Hard]}}.
The \textcolor{blue!60!black}{\textsf{[Hard]}} questions have several parts and ask the student to \emph{reason}, not to reproduce a
derivation from the lecture notes or the exercise sheets: a few careful paragraphs (with equations only where they
carry the argument) is the expected answer for each part.

\section{Euclidean path integrals and lattice scalar field theory}

\begin{question}
\diff{Easy} What is a Wick rotation, and why does rotating to Euclidean time make the path integral converge? Contrast the real-time weight $e^{iS/\hbar}$ with the Euclidean weight $e^{-S_E/\hbar}$.
\end{question}
\begin{solution}
A Wick rotation analytically continues the time variable $t \to -i\tau$ (equivalently rotating the time contour by $90^\circ$ in the complex plane). Under it the oscillatory real-time factor $e^{-iHt/\hbar}$ becomes the exponentially damped $e^{-H\tau/\hbar}$. In the path integral the weight changes from the pure phase $e^{iS/\hbar}$, whose rapid oscillations make the integral only conditionally convergent, to the real positive Boltzmann-like weight $e^{-S_E/\hbar}$ with $S_E \ge 0$. The latter is absolutely convergent and can be interpreted as a probability weight, which is what makes Monte Carlo sampling possible.
\end{solution}

\begin{question}
\diff{Intermediate} Write down the discretized Euclidean action for a free scalar field on a hypercubic lattice of spacing $a$, and explain how $a$ acts as an ultraviolet regulator via the Brillouin zone.
Both the choice $\Box_L = \Delta^b_\mu\Delta^f_\mu$ and the choice $\Box_L = \Delta^s_\mu\Delta^s_\mu$ (symmetric derivative squared) reduce to $\partial_\mu\partial_\mu$ as $a\to 0$. Using their momentum-space forms, explain why the second choice is nevertheless a bad discretization even for a \emph{scalar} field, and name the phenomenon in the fermionic case that it is the direct analogue of.
\end{question}
\begin{solution}
\begin{equation}
S_E = a^4 \sum_n \left[ \frac{1}{2}\sum_\mu \left(\frac{\phi(n+\hat\mu)-\phi(n)}{a}\right)^2 + \frac{1}{2} m^2 \phi(n)^2 \right].
\end{equation}
On a lattice of spacing $a$ the field lives only at discrete sites, so its Fourier transform is periodic with period $2\pi/a$ and momenta may be restricted to the first Brillouin zone, $p_\mu \in (-\pi/a, \pi/a]$. This provides a hard momentum cutoff $\Lambda \sim \pi/a$: no modes with wavelength shorter than $a$ exist. Loop integrals that would diverge in the continuum become finite integrals over a compact zone, so $a>0$ regulates the UV divergences.

For the second part, compare the momentum-space kernels. The forward--backward Laplacian gives
\begin{equation}
-\widetilde{\Box}(p) = \frac{4}{a^2}\sum_\mu \sin^2\!\left(\frac{p_\mu a}{2}\right),
\end{equation}
which vanishes \emph{only} at $p_\mu = 0$ inside the Brillouin zone. The symmetric derivative gives $\tilde\Delta^s_\mu(p) = \tfrac{i}{a}\sin(p_\mu a)$, hence
\begin{equation}
-\widetilde{\Box}^s(p) = \frac{1}{a^2}\sum_\mu \sin^2(p_\mu a),
\end{equation}
where the missing factor of one half in the argument is decisive: $\sin(p_\mu a)$ vanishes both at $p_\mu = 0$ \emph{and} at $p_\mu = \pm\pi/a$. The inverse propagator therefore has $2^d$ zeros in the Brillouin zone rather than one, and each of them behaves like a massless continuum pole. The lattice theory thus describes $2^d$ degenerate species instead of one, and the extra ones do not decouple as $a\to0$ because their contribution to momentum integrals is not suppressed. This is precisely the fermion doubling problem, which arises for naive lattice fermions for exactly the same reason: the naive Dirac operator is built from the \emph{symmetric} difference, and is linear in the derivative so no symmetric-derivative-squared alternative is available. For the scalar one simply avoids the problem by choosing $\Delta^b_\mu\Delta^f_\mu$; for the fermion the first-order structure of the Dirac operator makes the pathology unavoidable without further modification.
\end{solution}

\begin{question}
\diff{Hard} The lattice weight $e^{-S_E}$ is positive, which is usually summarized as ``a Euclidean lattice field theory is just a classical statistical system''. Consider the free scalar field on an $N_s^3\times N_t$ lattice with periodic boundary conditions, $S_E = \tfrac{1}{2}\sum_{n,m}\hat\phi_n K_{nm}\hat\phi_m$, and recall the identification $Z = \mathrm{Tr}\, e^{-\hat H/T}$ with $1/T = N_t a$.
\begin{enumerate}[label=(\roman*), leftmargin=*]
  \item Positivity of the Boltzmann weight is \emph{necessary but not sufficient} for the Euclidean lattice correlators to be interpretable as expectation values of operators in a Hilbert space. State the additional property the lattice action must have, and explain exactly what it buys you --- i.e.\ which statements about masses and about the spectrum would be unjustified without it.
  \item Explain why the zero-spatial-momentum correlator behaves as $\cosh\!\big[\hat E\,(\tau - N_t/2)\big]$ rather than as a single exponential $e^{-\hat E \tau}$. Identify precisely where the second exponential comes from in the trace formula, and state what happens to it as $N_t\to\infty$ at fixed $a$. Is this structure an artifact of the discretization, or unavoidable at any finite temperature?
  \item Suppose a lattice action produced a transfer matrix with a \emph{negative} eigenvalue $-|\lambda|$. What would you observe in the measured correlator $C(\tau)$, what ``energy'' would a naive exponential fit return, and what must happen to such a state as $a\to 0$ for the lattice theory to be acceptable?
\end{enumerate}
\end{question}
\begin{solution}
\textbf{(i)} What is needed in addition is that the action be local in Euclidean time and \emph{reflection positive} (Osterwalder--Schrader positivity): the weight must factorize across a time slice in such a way that the time evolution over one lattice step is generated by a \emph{transfer matrix} $\hat T$ which is Hermitian and positive definite, so that one may write $\hat T = e^{-a\hat H}$ with $\hat H$ Hermitian and bounded from below. Positivity of $e^{-S_E}$ only guarantees that the configuration weights can be read as probabilities of a classical statistical system --- it says nothing about whether that statistical system is the Euclidean version of a \emph{quantum} theory.

What reflection positivity buys is the spectral representation: it lets one write
\begin{equation}
C(\tau) = \langle \hat\phi(\tau)\hat\phi(0)\rangle = \frac{1}{Z}\,\mathrm{Tr}\!\left[\hat T^{\,N_t-\tau/a}\,\hat\phi\,\hat T^{\,\tau/a}\,\hat\phi\right]
= \sum_n |\langle 0|\hat\phi|n\rangle|^2 e^{-E_n \tau} + \dots
\end{equation}
with \emph{real, positive} coefficients and \emph{real} energies $E_n \geq E_0$. Without it, none of the standard lattice statements is justified: that the exponential decay rate of a correlator is a physical energy; that $E_0$ is the lowest energy in the channel so that the slowest exponential dominates at large $\tau$; that the overlaps are positive so a multi-exponential fit is a sum of positive contributions and effective masses approach the ground state monotonically from above; and that a Hilbert space with positive-definite norm exists at all. One could still have a perfectly well-behaved positive statistical system whose ``correlators'' decay with complex or oscillatory rates and have no interpretation as $|\langle 0|\hat\phi|n\rangle|^2 e^{-E_n\tau}$.

\textbf{(ii)} The trace in $Z = \mathrm{Tr}\,e^{-\hat H/T}$ is what produces the second exponential. Inserting the operator $\hat\phi$ at times $0$ and $\tau$ into the trace over the periodic temporal direction gives
\begin{equation}
C(\tau) = \frac{1}{Z}\sum_n |\langle 0|\hat\phi|n\rangle|^2 \left( e^{-E_n \tau} + e^{-E_n (1/T - \tau)} \right) + \dots
= \sum_n \frac{2|\langle 0|\hat\phi|n\rangle|^2}{Z} e^{-E_n/2T}\cosh\!\left[E_n\left(\tau - \frac{1}{2T}\right)\right].
\end{equation}
Physically: because time is periodic, the operator at separation $\tau$ is also at separation $1/T - \tau$ going the ``other way around'' the torus, and the state can propagate either way. Equivalently, the second term is the contribution in which the excitation propagates backwards in Euclidean time through the temporal boundary. In lattice units, with $\hat E = E a$ and $1/T = N_t a$, this is the quoted $\cosh[\hat E(\tau - N_t/2)]$, symmetric about $\tau = N_t/2$.

As $N_t \to \infty$ at fixed $a$ (i.e.\ $T\to0$), the backward term $e^{-E_n(1/T - \tau)}$ is suppressed by $e^{-E_n/T}$ and the correlator becomes a pure sum of decaying exponentials. This is \emph{not} a discretization artifact: it is a consequence of the trace, i.e.\ of the compactness of Euclidean time, and therefore occurs in the continuum theory at any finite temperature as well. It is only the finite \emph{temporal extent} that matters, not the finite spacing. (For the same reason a lattice calculation intended to extract vacuum masses must have $N_t a \gg 1/m$, and states with different temporal boundary conditions --- e.g.\ fermions with antiperiodic time --- produce a $\sinh$ instead of a $\cosh$.)

\textbf{(iii)} A negative transfer-matrix eigenvalue means $\hat T$ has an eigenvalue $-|\lambda| = -e^{-\hat E}$, so the corresponding contribution to the correlator is
\begin{equation}
(-|\lambda|)^{\tau/a} = (-1)^{\tau/a}\, e^{-\hat E \tau /a},
\end{equation}
i.e.\ an exponential decay multiplied by an alternating sign from time slice to time slice. Observationally one sees a correlator that zig-zags: even and odd time slices lie on two separate smooth curves, and the effective mass computed from consecutive slices oscillates instead of plateauing. A naive single-exponential fit returns a nonsensical value; the correct statement is that the state has $\hat T$ eigenvalue $-e^{-\hat E}$, i.e.\ an ``energy'' $E + i\pi/a$ --- the imaginary part being pure cutoff, of order the inverse lattice spacing. Such a state is therefore not a physical excitation but a lattice artifact (this is exactly how the parity partner appears in staggered-fermion correlators, where one deals with it by projecting onto even/odd slices or by fitting an oscillating ansatz). For the lattice theory to be acceptable, the artifact energy must diverge as $a\to 0$: $\pi/a \to \infty$, so the state decouples from the low-energy spectrum. If instead the negative-eigenvalue state survived at a fixed physical energy, the continuum limit would not be a unitary quantum field theory with a positive-definite Hilbert space.
\end{solution}

\section{Gauge theories on a lattice}

\begin{question}
\diff{Easy} Define the link variable $U_\mu(n)$ and state its gauge transformation law. Why is a link variable, rather than the field $A_\mu$ directly, the natural object on the lattice?
\end{question}
\begin{solution}
The link variable is the parallel transporter along the edge from site $n$ to $n+\hat\mu$,
\begin{equation}
U_\mu(n) = \exp\!\big(i a g A_\mu(n)\big) \in SU(N),
\end{equation}
and it transforms as $U_\mu(n) \to \Omega(n)\, U_\mu(n)\, \Omega(n+\hat\mu)^\dagger$. It is the natural variable because gauge invariance is realized \emph{exactly} on the lattice through group-valued links: products of links along closed loops are gauge invariant, and one never has to expand in the algebra-valued $A_\mu$. This keeps the lattice measure (the Haar measure over the compact group) gauge invariant to all orders in $a$, rather than only up to $\mathcal{O}(a)$ as a naive discretization of $D_\mu$ would be.
\end{solution}

\begin{question}
\diff{Intermediate} Define the plaquette $U_{\mu\nu}(n)$, write down the Wilson gauge action, and explain why the plaquette is the smallest gauge-invariant object built from links. Now confront two facts: (a) the Wilson action reduces to $\tfrac14\int d^4x\, F_{\mu\nu}F_{\mu\nu}$ as $a\to 0$, and (b) \emph{compact} lattice $U(1)$ has a strong-coupling phase in which static charges are confined, whereas continuum QED has no such phase. Explain how both can be true simultaneously, what feature of the lattice formulation is responsible, and what this implies about \emph{where} in coupling space the continuum limit of a lattice gauge theory may be taken.
\end{question}
\begin{solution}
The plaquette is the ordered product of links around an elementary square,
\begin{equation}
U_{\mu\nu}(n) = U_\mu(n)\, U_\nu(n+\hat\mu)\, U_\mu^\dagger(n+\hat\nu)\, U_\nu^\dagger(n),
\end{equation}
and the Wilson gauge action is
\begin{equation}
S_W = \beta \sum_n \sum_{\mu<\nu}\left(1 - \frac{1}{N}\,\mathrm{Re}\,\mathrm{Tr}\, U_{\mu\nu}(n)\right), \qquad \beta = \frac{2N}{g^2}
\end{equation}
(for $U(1)$, $S_G = \frac{1}{e^2}\sum_n\sum_{\mu<\nu}\mathrm{Re}\,[1 - U_{\mu\nu}(n)]$). A single link is not gauge invariant, since it transforms with $\Omega(n)$ at one end and $\Omega(n+\hat\mu)^\dagger$ at the other. Only a \emph{closed} loop has all intermediate gauge factors cancel; the plaquette is the shortest closed loop on a hypercubic lattice, hence the smallest gauge-invariant object.

Both statements (a) and (b) are true because they concern different regions of coupling space, and the $a\to0$ expansion is not a statement about the whole theory. The expansion $U_{\mu\nu} = \exp(ia^2 e F_{\mu\nu} + \dots)$ and the subsequent identification with $\tfrac14 F^2$ assume the plaquette angle $\theta_{\mu\nu}$ is \emph{small}, i.e.\ typical field fluctuations are weak. This is a weak-coupling (large-$\beta$) statement. The feature responsible for the discrepancy is \textbf{compactness}: the lattice variable is the group element $U_\mu = e^{i\theta_\mu}$ with $\theta_\mu \in (-\pi,\pi]$, so the action is a \emph{periodic} function of the plaquette angle, whereas the continuum action $\tfrac14 F^2$ is a single quadratic well. Periodicity admits configurations with no continuum counterpart --- in particular Dirac-string/monopole excitations, where the plaquette angle winds by $2\pi$ around a closed surface. At strong coupling (small $\beta$) these topological excitations are abundant and condense, disordering the Wilson loop and producing an area law: static charges are confined. This confining phase is a property of the \emph{compact lattice} model, not of continuum QED; it is separated from the weak-coupling Coulomb phase by a phase transition at some $\beta_c$.

The implication is important and general: the continuum limit may only be taken in the phase whose long-distance physics is the intended continuum theory, approaching the appropriate critical point from the correct side. For compact $U(1)$ this is the weak-coupling Coulomb phase ($\beta > \beta_c$, where $a\to0$ requires $e\to0$); the strong-coupling confining phase is a lattice artifact phase with no continuum interpretation. More broadly: a lattice action that reproduces the right classical continuum limit under a small-$a$ expansion is not thereby guaranteed to have the right nonperturbative phase structure everywhere --- one must check which phase one is simulating in.
\end{solution}

\begin{question}
\diff{Hard} On the lattice one does not need to fix a gauge, but one may: the standard construction sets the links of a \emph{maximal tree} to the identity.
\begin{enumerate}[label=(\roman*), leftmargin=*]
  \item On a $D$-dimensional periodic lattice with $N$ sites there are $DN$ links. How many of them can be set to $\mathbb{1}$ by a gauge transformation, and --- crucially --- why not more? How many links are left to integrate over, and how does that number compare with the naive count ``$DN$ links minus one gauge parameter per site''?
  \item A Polyakov loop is a product of temporal links winding once around the periodic time direction. Show that no choice of maximal tree can set \emph{all} of its links to $\mathbb{1}$, and state the structural property that distinguishes the loops which can be trivialized from those which cannot. What does this say about which degrees of freedom survive complete gauge fixing?
  \item Elitzur's theorem states that $\langle U_\mu(n)\rangle = 0$ exactly, with no gauge fixing at all. Yet in the maximal-tree gauge $U_\mu(n)=\mathbb{1}$ on every tree link, so its ``expectation value'' is $1$. Resolve the apparent contradiction. Then state precisely which class of quantities is guaranteed to be independent of the gauge choice, and give the one-line argument for why $\langle U_\mu(n)\rangle$ must vanish in the unfixed theory.
\end{enumerate}
\end{question}
\begin{solution}
\textbf{(i)} There is one gauge matrix $\Omega(n)$ per site, so $N$ parameters (matrices), but a \emph{global} transformation --- the same $\Omega$ everywhere --- does not help: for $U(1)$ it leaves every link untouched, and for a non-abelian group it merely conjugates the whole configuration, so it cannot trivialize an additional link. Hence at most $N-1$ links can be set to the identity, and the set of trivialized links must form a \emph{tree}: a connected subgraph containing no closed loop. The reason is the mechanism seen in the lecture --- to set $U_\mu(n)\to\mathbb{1}$ one uses up the gauge freedom at one of its two endpoints, and each site's freedom can only be spent once. If one tried to trivialize a set of links containing a closed loop, one would be trying to fix more links than there are sites in the loop, and the last link would be determined by the others: the ordered product around a closed loop is gauge invariant and cannot be changed at all (see (ii)). A maximal (spanning) tree on $N$ sites has exactly $N-1$ links, so
\begin{equation}
\#\{\text{links set to }\mathbb{1}\} = N-1, \qquad \#\{\text{remaining links}\} = DN - (N-1) = (D-1)N + 1 .
\end{equation}
For $D=4$ this leaves $3N+1$ links instead of $4N$. The comparison with the naive count is the point of the question: ``$DN$ minus one parameter per site'' would give $DN - N = (D-1)N$, i.e.\ one link \emph{fewer} than the correct answer. The discrepancy of exactly $+1$ is the global gauge transformation, which is a redundancy of the redundancy: it acts trivially on the links, so it cannot be used to remove a degree of freedom. (This is the lattice analogue of the statement that the gauge group acting effectively on fields is $\mathcal{G}/\mathcal{G}_{\rm global}$.)

\textbf{(ii)} Let $P(n) = \prod_{n_t=0}^{N_t-1} U_4(n,n_t)$ be the Polyakov loop. Under a gauge transformation the intermediate factors telescope,
\begin{equation}
P(n) \to \Omega(n,0)\,P(n)\,\Omega(n,N_t)^\dagger = \Omega(n,0)\,P(n)\,\Omega(n,0)^\dagger ,
\end{equation}
using periodicity $\Omega(n,N_t) = \Omega(n,0)$. Hence $\mathrm{Tr}\,P(n)$ is strictly gauge invariant: it is a number that no gauge transformation can alter. If all $N_t$ links of the loop could be set to $\mathbb{1}$, then $\mathrm{Tr}\,P$ would equal $N$ (the dimension of the representation) on every configuration, contradicting gauge invariance for any configuration where it does not. Structurally: a maximal tree contains \emph{no closed loops by construction}, so a winding loop necessarily contains at least one non-tree link; and this is forced --- the loop is \emph{non-contractible}, it closes only because of the periodic boundary conditions and cannot be shrunk to a point, so it is not the boundary of any set of plaquettes whose links could all be absorbed into a tree.

The general statement: gauge fixing removes only the redundancy, never physics. What survives complete gauge fixing is precisely the gauge-invariant content --- the holonomies. Contractible loops (plaquettes and their products) survive as products of the leftover non-tree links; and in addition there are the $D$ independent non-contractible windings around the torus, whose traces are gauge invariant on their own. This is exactly why the Polyakov loop is a legitimate order parameter for deconfinement while the link itself is not.

\textbf{(iii)} There is no contradiction, because the two ``expectation values'' are of different objects. Elitzur's theorem is a statement about the \emph{unfixed} path integral, in which one sums over all gauge-equivalent copies of every configuration. In the maximal-tree gauge one has \emph{restricted} the integration domain to one representative per gauge orbit, chosen precisely such that the tree links are $\mathbb{1}$. The number ``$1$'' obtained there is not the same functional integral; it is a property of the gauge slice one selected, and it would change if a different maximal tree (or a different gauge condition) were chosen --- recall that the maximal tree is not unique. A gauge-variant operator simply does not have a gauge-independent expectation value; asking for $\langle U_\mu(n)\rangle$ is asking a question with no gauge-invariant answer, and the two computations return two different, equally meaningless answers.

What \emph{is} guaranteed to be gauge independent is the expectation value of any gauge-invariant functional $O[U]$ --- traces of closed loops (Wilson loops, plaquettes, Polyakov loops) and their products, and anything built from them. Since the action and the Haar measure are gauge invariant, restricting to a gauge slice changes both numerator and denominator by the same volume factor of the gauge orbit and leaves $\langle O\rangle$ unchanged. This is precisely why gauge fixing on the lattice is optional: it is done for convenience (fewer integration variables, or to make an interpretation such as the Wilson-loop-as-two-Wilson-lines transparent in temporal gauge), never out of necessity.

The one-line argument for Elitzur: perform, at fixed configuration, the gauge transformation at the single site $n$ only. Then $U_\mu(n) \to \Omega(n)U_\mu(n)$ while both the action and the measure are unchanged, so
\begin{equation}
\langle U_\mu(n)\rangle = \int d\Omega(n)\,\Omega(n)\,\langle U_\mu(n)\rangle = 0,
\end{equation}
since the Haar average of a group element over a compact group vanishes ($\int_{-\pi}^{\pi} \frac{d\theta}{2\pi} e^{i\theta} = 0$ for $U(1)$; $\int d\Omega\,\Omega = 0$ for $SU(N)$). The averaging over the \emph{local} gauge freedom, which is present configuration by configuration and is not broken by any finite-volume boundary condition, kills the expectation value exactly --- no spontaneous breaking of a local symmetry is possible.
\end{solution}

\section{Fermions on a lattice}

\begin{question}
\diff{Easy} State the fermion doubling problem. How many doublers does the naive discretization of the Dirac operator produce in $d$ dimensions, and where in the Brillouin zone do they sit?
\end{question}
\begin{solution}
Replacing $\gamma_\mu\partial_\mu$ by the symmetric difference $\gamma_\mu\frac{\psi(n+\hat\mu)-\psi(n-\hat\mu)}{2a}$ gives an inverse propagator $\propto \frac{i}{a}\sum_\mu \gamma_\mu \sin(p_\mu a)$, which vanishes not only at $p_\mu=0$ but at every $p_\mu \in \{0, \pi/a\}$. Each momentum component thus contributes two zeros, giving $2^d$ degenerate species; in $d=4$ this is $16$ fermions where the theory should describe one. The extra $15$ doublers sit at the corners and edge-centres of the Brillouin zone (momentum components equal to $\pi/a$).
\end{solution}

\begin{question}
\diff{Intermediate} Compare Wilson and staggered fermions: by what mechanism does each remove (or reduce) the doublers, and what is the price? In the lecture's statement of the no-go theorem, Wilson fermions were said to break \emph{chiral symmetry} while staggered fermions break \emph{translation invariance}. But staggered fermions are also routinely described as retaining a remnant chiral symmetry, and they are manifestly invariant under translations by $2a$. Reconcile these statements: what exactly is broken, at which scale, and why does the theorem not care that a larger-period translation invariance survives?
\end{question}
\begin{solution}
\emph{Wilson fermions} add a dimension-five term $-\frac{r}{2}\sum_n \bar\psi\,\hat\Box\,\psi$, which makes the mass momentum dependent, $M(p) = M + \frac{2r}{a}\sum_\mu \sin^2(p_\mu a/2)$. At $p=0$ this is $M$, but at each Brillouin-zone corner it is enhanced by $2r/a$ per component, so the $15$ doublers acquire masses of order $1/a$ and decouple as $a\to0$. The price is that the Wilson term explicitly breaks chiral symmetry at any finite $a$, even for $M=0$, which forces an additive mass renormalization (fine-tuning to the critical mass) and allows $\mathcal{O}(a)$ discretization errors. \emph{Staggered fermions} instead distribute the spinor components over the sites of an elementary hypercube (spin diagonalization), so that the effective spacing for a given component is $2a$ and the effective Brillouin zone is halved. This reduces $2^d$ species to $2^{d/2}$ ``tastes'' ($16 \to 4$ in $d=4$), preserves a remnant $U(1)$ chiral symmetry that protects the mass from additive renormalization, but leaves a taste degeneracy that must be dealt with (e.g.\ the fourth-root prescription) and generates taste-breaking artifacts.

The reconciliation: what the staggering transformation breaks is invariance under translations by a \emph{single} lattice spacing. The transformation $T(n) = (\gamma^0\gamma^1)^{n_1}\cdots(\gamma^0\gamma^d)^{n_d}$ and the resulting sign factors $\eta_j(n) = (-1)^{n_1+\dots+n_{j-1}}$ depend on the site \emph{parity}, so the action is not invariant under $n \to n+\hat\mu$; it is only invariant under $n \to n + 2\hat\mu$, together with the accompanying shift symmetry that mixes tastes. So the correct statement is: the symmetry group of one-site translations is broken down to a subgroup of two-site translations. That is precisely the assumption of the Nielsen--Ninomiya theorem that fails, because ``translation invariance'' there means invariance under the \emph{full} lattice translation group with the fundamental spacing $a$ --- that is what makes the momentum-space analysis on the Brillouin zone $(-\pi/a,\pi/a]$ apply, and what forbids the operator from distinguishing sublattices. Once one only has $2a$ translations, the relevant zone is $(-\pi/2a,\pi/2a]$ per direction and the counting of zeros changes: the doubling is halved per direction, which is exactly the $2^d \to 2^{d/2}$ reduction.

Similarly, the ``remnant chiral symmetry'' is not the full chiral symmetry either: the staggered operator does not satisfy $\{D,\gamma_5\}=0$ with the continuum $\gamma_5$ acting on a four-component spinor at a single site --- there is no such spinor. What survives is a single $U(1)$ generated by the site-parity factor $\tilde\varepsilon(n) = (-1)^{\sum_i n_i}$, a subgroup of the continuum $U(4)\times U(4)$ taste-chiral symmetry. So neither statement is wrong; both symmetries are reduced to subgroups, and the theorem's assumptions refer to the full symmetries. The general lesson is that the no-go theorem is evaded by weakening an assumption to a subgroup or to an $\mathcal{O}(a)$-violated relation (as the Ginsparg--Wilson relation $\{D,\gamma_5\} = aD\gamma_5 D$ does for chirality), not by satisfying all four exactly.
\end{solution}

\begin{question}
\diff{Hard} Rather than restating the Nielsen--Ninomiya theorem, put its four assumptions (locality, translation invariance, $\{D,\gamma_5\}=0$, correct continuum limit) to work.
\begin{enumerate}[label=(\roman*), leftmargin=*]
  \item The naive lattice Dirac operator satisfies locality, translation invariance and exact chiral symmetry. Which assumption does it violate, and how is that visible in the momentum-space operator? Be precise about \emph{why} the failure is a failure --- what goes wrong in the continuum limit of a correlator, as opposed to merely ``there are extra poles''.
  \item Now define the Dirac operator directly in momentum space by $\tilde D(p) = i\sum_\mu \gamma_\mu p_\mu$ with $p_\mu$ restricted to the first Brillouin zone (the ``SLAC'' derivative). Which of the four assumptions fails now, and how does the failure show up in position space? Explain why it is fatal in an \emph{interacting} theory, rather than merely inelegant --- what does the offending assumption actually protect?
  \item The $2^d$ species of the naive operator do not all carry the same chirality. Use this to explain why naive lattice fermions, despite possessing an exact chiral symmetry, cannot reproduce the axial anomaly of the continuum theory. Then argue why this makes the no-go theorem look inevitable rather than accidental, and identify what supplies the anomaly in the Wilson formulation.
\end{enumerate}
\end{question}
\begin{solution}
\textbf{(i)} The naive operator violates assumption (iv), the correct continuum limit: it does not describe a single Dirac fermion. In momentum space
\begin{equation}
\tilde D_{\rm naive}(p) = \frac{i}{a}\sum_\mu \gamma_\mu \sin(p_\mu a) + m,
\end{equation}
and for $m=0$ the operator has $2^d$ zeros in the Brillouin zone, at $p_\mu \in \{0,\pi/a\}$, because $\sin(p_\mu a)$ --- unlike the scalar's $\sin(p_\mu a/2)$ --- vanishes at the zone boundary as well.

The precise reason this is a failure, rather than a curiosity, is what happens in a momentum integral. Consider the propagator inserted into a two-point function, $\int_{\rm BZ} d^dp\, \tilde D^{-1}(p)\,e^{ipx}$. In taking $a\to0$ one wants to argue that only modes with $|p| \ll 1/a$ contribute, so that $\frac{1}{a}\sin(p_\mu a) \to p_\mu$ may be substituted and the continuum expression recovered. That argument works for the scalar because the inverse propagator is large ($\sim 1/a^2$) everywhere except near $p=0$, so the zone-boundary region is suppressed. Here it fails: near \emph{each} of the $2^d$ zeros the inverse propagator is again small and linear in the deviation, so each region contributes to the integral with the same weight as $p\approx0$ and survives the limit. In the language of the lecture, the modes at the corners of the Brillouin zone contribute equally in the continuum limit. Every correlator therefore comes out $2^d$ times too large (more precisely, as that of a theory of $2^d$ degenerate species): the extra states are physical excitations of the lattice theory at arbitrarily low energy, not high-energy artifacts, and they do not decouple as the cutoff is removed. That is what distinguishes doublers from harmless lattice artifacts.

\textbf{(ii)} Now locality (i) fails. The operator is translation invariant, it manifestly anticommutes with $\gamma_5$, and it has a single zero at $p=0$ with exactly the continuum dispersion inside the whole zone --- so (ii), (iii) and (iv) hold at the free level. But $\tilde D(p) = i\gamma_\mu p_\mu$ restricted to $p_\mu\in(-\pi/a,\pi/a]$ and continued periodically is a \emph{sawtooth}: it is discontinuous at the zone boundary. A Fourier coefficient of a discontinuous function falls off only as a power, so in position space the operator is a hopping term connecting \emph{all} pairs of sites with an amplitude decaying like $1/|n-m|$ rather than exponentially. There is no finite range: it is a non-local action.

Why fatal? Locality is not an aesthetic requirement; it is what guarantees that the lattice theory has a well-defined local continuum limit at all. Specifically, locality (exponentially bounded couplings) is what ensures that the effective action obtained by integrating out the cutoff scale is a sum of \emph{local} operators of increasing dimension --- the Symanzik expansion --- so that the artifacts appear as powers of $a$ times local operators and vanish as $a\to0$, and that the theory is renormalizable with a finite number of counterterms. It is also what guarantees that the continuum limit respects Euclidean (Lorentz) invariance and cluster decomposition. With the SLAC derivative none of that holds: the discontinuity at the zone boundary produces, in loop diagrams of the interacting theory, contributions from the zone boundary that survive $a\to0$ and generate non-covariant, non-local terms in the effective action --- the Karsten--Smit analysis shows that the correct covariant continuum Green's functions are not recovered, and the required counterterms are non-local. In the free theory nothing is wrong, which is precisely why one must state the assumption for the interacting theory. So the SLAC operator does not evade the no-go theorem: it exhibits it, by showing what the price of keeping chirality and losing locality actually is.

\textbf{(iii)} Expand the naive operator around each of its $2^d$ zeros, $p = \bar p_A + k$ with $\bar p_A$ a corner of the zone having some subset of components equal to $\pi/a$. Near each zero one finds a continuum-like Dirac operator, but with $\gamma$-matrices rotated by a unitary transformation $\Gamma_A$ built from products of $\gamma_\mu$ over the directions in which $\bar p_A$ has a component $\pi/a$ (because $\sin(\pi/a + k)a = -\sin(ka)$ flips the sign of those terms). The relative chirality of species $A$ is then $\pm1$ according to whether that number of components is even or odd, so of the $2^d$ species exactly half have chirality $+$ and half have chirality $-$. In $d=4$: $8$ and $8$.

Consequently the naive lattice theory describes a \emph{vector-like, anomaly-free} set of species: the would-be axial anomaly of each species is cancelled pairwise by a partner of opposite chirality. In the continuum the axial anomaly is a genuine physical prediction (the non-conservation of the singlet axial current, $\partial_\mu j^\mu_5 \propto \mathrm{Tr}\,F\tilde F$, which is measured in $\pi^0 \to \gamma\gamma$ and controls the $\eta'$ mass). A lattice theory with an exact chiral symmetry has a chirally invariant measure and action, hence an exactly conserved axial current and no anomaly at all --- so it \emph{cannot} be the regularization of a single continuum Dirac fermion. The doublers are precisely the mechanism by which the lattice reconciles exact chirality with the absence of an anomaly: the anomaly does not vanish species by species, it cancels between species.

This is why the theorem looks inevitable rather than accidental. Turn the argument around: suppose an operator existed that was local, translation invariant, exactly chirally symmetric \emph{and} described one species. Then the lattice theory would have an exactly conserved singlet axial current with a single fermion --- contradicting a robust, verified feature of the continuum theory it is supposed to regulate. So the obstruction is not a technical accident of the difference-operator construction; it is the statement that a chirally symmetric regulator cannot produce an anomaly, and something must give. (This is also why the Ginsparg--Wilson relation $\{D,\gamma_5\} = aD\gamma_5D$ is exactly the right amount of violation: the modified symmetry is exact, but the fermion \emph{measure} is not invariant under the L\"uscher transformation, and its Jacobian reproduces the anomaly and the index theorem.)

In the Wilson formulation, the Wilson term is what supplies the anomaly. It breaks chiral symmetry explicitly by $\mathcal{O}(a)$, and as the doublers are given masses $\sim 1/a$ and removed, the cancellation between opposite-chirality species is spoiled; in the $a\to0$ limit the chirality-violating Wilson vertex leaves behind precisely the finite, correct continuum anomaly. The lesson is that the chiral-symmetry breaking of Wilson fermions is not merely a cost to be tolerated --- part of it is physically required.
\end{solution}

\section{Continuum limit}

\begin{question}
\diff{Easy} What is the physical correlation length $\xi_{\rm phys}$ in terms of the lattice spacing $a$ and the dimensionless correlation length $\xi$? Why must $\xi \to \infty$ (in lattice units) to reach the continuum?
\end{question}
\begin{solution}
The physical correlation length is $\xi_{\rm phys} = a\,\xi$, where $\xi$ is measured in lattice units (the number of lattice sites over which correlations decay). Equivalently $\xi = 1/(a m_{\rm phys})$ for a physical mass $m_{\rm phys}$. To take $a\to 0$ while keeping $\xi_{\rm phys}$ (a physical quantity) fixed, the dimensionless $\xi = \xi_{\rm phys}/a$ must diverge. A diverging correlation length in lattice units is precisely the hallmark of a critical point, so the continuum limit lives at a critical point of the lattice model.
\end{solution}

\begin{question}
\diff{Intermediate} Explain why the continuum limit must be taken at a second-order (critical) point of the lattice theory, and why a first-order transition is unsuitable. Then resolve the following: the deconfinement transition of pure $SU(3)$ gauge theory at finite temperature \emph{is} first order, and it is nevertheless studied on the lattice with a perfectly well-defined continuum limit. Why is this not a contradiction? Distinguish carefully the two transitions involved and the parameter that is tuned in each.
\end{question}
\begin{solution}
At a second-order transition the correlation length diverges \emph{continuously} and there is no latent heat, so one can tune the bare couplings to make $\xi \to \infty$ smoothly; this is what lets $a = \xi_{\rm phys}/\xi \to 0$ while physics stays fixed, and it also guarantees the loss of memory of the lattice details (universality). At a first-order transition the correlation length stays \emph{finite} across the transition --- two phases coexist with a discontinuity and latent heat --- so $\xi$ never diverges, $a$ cannot be sent to zero, and there is no scaling region in which lattice artifacts vanish.

The resolution of the apparent paradox is that two entirely different transitions are involved, tuned by different parameters and living at different scales.
\begin{itemize}
  \item The transition relevant to the \emph{continuum limit} is a critical point in the space of \emph{bare lattice couplings} at zero physical temperature, reached by tuning $\beta \to \infty$ (i.e.\ $g\to0$) so that $a\to0$: the Gaussian UV fixed point of the asymptotically free theory. It is here that $\xi$ measured \emph{in lattice units} must diverge. For $SU(3)$ this fixed point exists and is approached as $g\to0$; that is what makes a continuum limit possible at all.
  \item The deconfinement transition is a \emph{thermal} transition of the already-continuum theory, located at a physical temperature $T_c \approx 270$~MeV, i.e.\ at a specific value of the physical quantity $T = 1/(N_t a)$. It is first order in pure $SU(3)$, meaning the physical correlation length $\xi_{\rm phys}$ stays finite there --- but $\xi_{\rm phys}$ finite is no obstacle whatsoever to $a \to 0$, since one only needs $\xi = \xi_{\rm phys}/a \to \infty$.
\end{itemize}
Concretely, one studies the thermal transition at fixed $N_t$ by tuning $\beta$ to bring $T = 1/(N_t a(\beta))$ to $T_c$, and then takes the continuum limit by repeating this at $N_t = 6, 8, 10, \dots$, each time with a larger $\beta$ (smaller $a$), and extrapolating $a\to0$. The dimensionless $\xi$ at the transition then grows like $\xi_{\rm phys}/a$ even though $\xi_{\rm phys}$ itself is finite and order-of-transition-independent. So the statement ``the continuum limit requires a second-order transition'' refers specifically to the critical behaviour in the bare parameters as $a\to0$, not to the order of any physical phase transition the theory happens to possess. A first-order \emph{bare-coupling} transition --- such as the strong-coupling bulk transitions of some lattice actions --- \emph{is} an obstruction, and one must keep the simulation away from it.
\end{solution}

\begin{question}
\diff{Hard} The continuum limit is a statement about renormalization-group flow, not about a formula for $a(g)$.
\begin{enumerate}[label=(\roman*), leftmargin=*]
  \item Two lattice actions --- the Wilson plaquette action, and an action that additionally includes $2\times 1$ rectangular loops --- give \emph{different} numbers for the same dimensionless ratio at the same lattice spacing. Explain in RG language why they must nonetheless agree in the continuum limit. What roles do the critical surface, the relevant directions and the irrelevant directions play, and what does universality \emph{not} guarantee?
  \item A student runs simulations at steadily increasing $\beta$ on a fixed $32^4$ lattice and reports that they are approaching the continuum limit. Identify the flaw, state what must be held fixed instead, and express the requirement as a condition relating $\xi$, $N$ and the physical size of the box. Name the two practical costs of doing it correctly, and say which lattice observable would first betray the student's error.
  \item The pure-gauge Wilson action has leading discretization errors of $\mathcal{O}(a^2)$, whereas Wilson fermions introduce $\mathcal{O}(a)$ errors. Explain this difference from the symmetries of the lattice action and the dimensions of the operators allowed in a Symanzik effective-action analysis --- in particular, why the offending fermionic operator is permitted for Wilson fermions but forbidden for a Ginsparg--Wilson operator. Finally, explain why improvement cannot make discretization errors vanish at finite $a$, only change which power of $a$ dominates.
\end{enumerate}
\end{question}
\begin{solution}
\textbf{(i)} In RG language, a lattice action is a point in the infinite-dimensional space of couplings $\{g_i\}$, and adding $2\times1$ rectangles simply moves that point: the two actions are two different starting points. Both, however, lie on (or are tuned onto) the \emph{critical surface} --- the basin of attraction of the same UV fixed point, here the Gaussian fixed point at $g=0$ of the asymptotically free theory. Under repeated RG (blocking) transformations, the flow from either starting point is attracted along the same unstable manifold: the components along the \emph{irrelevant} directions shrink to zero, while the \emph{relevant} directions (in pure gauge, just the single marginally relevant coupling $g$, plus the quark masses when fermions are present) must be tuned --- this is what one does when choosing $\beta$ to fix $a$. Since the two actions differ only by irrelevant operators --- the rectangle term is, in the Symanzik language, the continuum action plus higher-dimension operators with different coefficients --- their flows converge, and all dimensionless ratios of physical quantities agree in the limit. Their disagreement at fixed finite $a$ is exactly the size of the irrelevant-operator contributions, i.e.\ the lattice artifacts, and one may legitimately \emph{use} that disagreement: agreement of continuum extrapolations from two different discretizations is the standard nontrivial check of a lattice result, and choosing the coefficients of the extra loops to cancel the leading artifact is the Symanzik-improvement programme.

Universality does \emph{not} guarantee: that the two actions give the same answers at finite $a$ (they do not, and there is no sense in which one is ``more right'' at finite $a$); that they approach the limit with the same rate or even the same leading power of $a$ (that depends on which irrelevant operators are present); that the approach is monotonic, or that either is close to the limit at the $a$ one can afford; nor that two actions belong to the same universality class in the first place --- that requires that they share the same symmetries and field content and lie on the same critical surface. A lattice action with a different symmetry (say, one that breaks a chiral symmetry, or that sits in a different phase) will happily flow somewhere else.

\textbf{(ii)} The flaw is that increasing $\beta$ at fixed $N$ decreases $a$ and therefore \emph{shrinks the physical volume}: $L_{\rm phys} = N a(\beta) \to 0$. The student is not approaching the continuum limit but squeezing the theory into an ever-smaller box, which is a completely different limit (and eventually a different physical regime --- at small volume the theory is weakly coupled and, for gauge theories, deconfined). What must be held fixed is the physical box size, so that $N$ must grow as $a$ shrinks:
\begin{equation}
L_{\rm phys} = N a = \text{const}, \qquad N = \frac{L_{\rm phys}}{a} \;\propto\; \frac{1}{a},
\end{equation}
and to control finite-size effects the box must also be large compared with the longest physical correlation length,
\begin{equation}
\xi \;=\; \frac{\xi_{\rm phys}}{a} \;=\; \frac{1}{a m_{\rm phys}} \;\ll\; N,
\qquad\text{i.e.}\qquad a \;\ll\; \xi_{\rm phys} \;\ll\; N a .
\end{equation}
This is the double requirement: $a$ small enough that discretization errors are small (many lattice spacings per correlation length), and $Na$ large enough that finite-volume errors are small (many correlation lengths per box). In practice one demands $m_\pi L \gtrsim 4$.

The two practical costs are (a) sheer computational cost: keeping $L_{\rm phys}$ fixed means $N \propto 1/a$, so the number of degrees of freedom grows like $a^{-4}$ and the total cost like a high power of $1/a$ (worse still once the algorithmic cost of the fermion solve is included); and (b) \emph{critical slowing down}: since one is deliberately approaching a critical point, $\xi \to \infty$ in lattice units, and Monte Carlo autocorrelation times grow as a power of $\xi$ --- for gauge theories, topological charge freezing is the notorious extreme case --- so each independent configuration costs more as well.

The observable that betrays the error first is any quantity sensitive to the box: the lowest-lying hadron mass in lattice units, $\hat m = m_{\rm phys}a$, decreases with $\beta$ as it should, but the \emph{dimensionless} combination $\hat m N = m_{\rm phys} L_{\rm phys}$ falls, and mass ratios start to drift away from their physical values as $m_\pi L$ drops below a few. For a gauge theory, the Polyakov loop is the cleanest tell: at fixed $N_t$, shrinking $a$ raises $T = 1/(N_t a)$, so the Polyakov loop will suddenly develop a nonzero expectation value --- the student's ``continuum limit'' has walked the simulation into the deconfined phase.

\textbf{(iii)} In the Symanzik framework, the lattice theory at small $a$ is described by an effective continuum action
\begin{equation}
S_{\rm eff} = S_{\rm cont} + a\sum_i c_i^{(5)} \mathcal{O}_i^{(5)} + a^2 \sum_i c_i^{(6)} \mathcal{O}_i^{(6)} + \dots,
\end{equation}
where the operators $\mathcal{O}_i^{(n)}$ of dimension $n$ must respect \emph{all} the symmetries of the lattice action: gauge invariance, the hypercubic subgroup of rotations, parity, charge conjugation, and any exact internal symmetry.

For the pure-gauge Wilson action, the available symmetries eliminate every dimension-five candidate. A dimension-five gauge-invariant operator would have to be of the form $\mathrm{Tr}\,F_{\mu\nu}D_\mu F_{\rho\nu}$-type with an odd number of derivatives/field strengths, and such structures either vanish identically by the Bianchi identity and the equations of motion, are total derivatives (irrelevant for the action), or are forbidden by parity and hypercubic symmetry. The first surviving corrections are dimension six ($\mathrm{Tr}\,D_\mu F_{\nu\rho}D_\mu F_{\nu\rho}$, $\mathrm{Tr}\,D_\mu F_{\mu\nu}D_\rho F_{\rho\nu}$, and the hypercubic-symmetric-but-not-$O(4)$-symmetric $\sum_\mu \mathrm{Tr}\,D_\mu F_{\mu\nu}D_\mu F_{\mu\nu}$), whence $\mathcal{O}(a^2)$. Equivalently and more simply: the Wilson plaquette action is even under $a \to -a$ in the sense that the plaquette is real and its expansion contains only even powers of $a$.

With Wilson fermions there \emph{is} an allowed dimension-five operator: the clover term
\begin{equation}
\mathcal{O}_{\rm cl} = \bar\psi\,\sigma_{\mu\nu}F_{\mu\nu}\,\psi ,
\end{equation}
together with $m\,\bar\psi\psi$ and $m^2\bar\psi\psi$-type terms that are absorbed into mass and coupling renormalization. The clover operator is gauge invariant, hypercubically invariant, parity and $C$ even --- but it is \emph{chirality flipping}: it connects left- and right-handed components, so it is forbidden by chiral symmetry. The Wilson term breaks chiral symmetry explicitly, hence nothing forbids it and its coefficient is nonzero: leading errors are $\mathcal{O}(a)$. This is also the cure --- adding the clover term to the action with a tuned coefficient $c_{\rm SW}$ cancels this contribution, which is exactly Sheikholeslami--Wohlert improvement. For a Ginsparg--Wilson operator (overlap, domain wall in the appropriate limit) the exact lattice chiral symmetry forbids $\bar\psi\sigma_{\mu\nu}F_{\mu\nu}\psi$ outright, so no tuning is needed and the leading errors are automatically $\mathcal{O}(a^2)$. Note the general moral: which power of $a$ one gets is dictated by the \emph{symmetries} one has chosen to keep, and the same reasoning explains why staggered fermions --- with their reduced translation and taste symmetry --- allow taste-violating operators that appear as their own class of $\mathcal{O}(a^2)$ artifacts.

Finally, improvement cannot remove artifacts entirely at finite $a$, for three independent reasons. First, the Symanzik expansion is an infinite tower: cancelling the $\mathcal{O}(a)$ (or $\mathcal{O}(a^2)$) term leaves the next order untouched, and full removal to all orders would require tuning infinitely many operator coefficients --- a perfect action. Second, improvement of the \emph{action} does not improve the \emph{operators} whose correlators one measures; those need their own improvement coefficients. Third, the coefficients themselves ($c_{\rm SW}$ and friends) are only known to finite accuracy, whether perturbatively (leaving a residual of the cancelled order times a power of $g^2$) or nonperturbatively (with a statistical/systematic error). Improvement therefore changes which power of $a$ dominates the approach to the continuum --- turning an $\mathcal{O}(a)$ extrapolation into an $\mathcal{O}(a^2)$ one, which is a large practical gain --- but a continuum extrapolation over several lattice spacings remains mandatory.
\end{solution}

\section{Chemical potential and the numerical sign problem}

\begin{question}
\diff{Easy} How is a quark chemical potential $\mu$ introduced into the lattice action, and why can it not simply multiply the number operator naively? State qualitatively why $\mu \neq 0$ leads to a complex action.
\end{question}
\begin{solution}
A chemical potential couples to conserved quark number, and one is tempted to add $\mu\,\bar\psi\gamma_4\psi$ to the Lagrangian. On the lattice this naive term produces a $\mu^2/a^2$ divergence in the energy density as $a\to0$ (Hasenfratz--Karsch). The cure is to introduce $\mu$ exactly as the fourth component of a constant imaginary abelian gauge field, $A_4 = -i\mu$, i.e.\ by multiplying the forward and backward temporal hopping terms by $e^{+\mu a}$ and $e^{-\mu a}$ respectively. This exponential (link-like) coupling is finite in the continuum limit and reproduces the correct density term at $\mathcal{O}(a)$. Because the two temporal directions are now weighted asymmetrically ($e^{+\mu a}\neq e^{-\mu a}$ for real $\mu$), the fermion matrix loses its $\gamma_5$-Hermiticity and its determinant becomes complex.
\end{solution}

\begin{question}
\diff{Intermediate} Explain the origin of the sign problem in terms of $\gamma_5$-Hermiticity of the Dirac operator, and why reweighting from a positive ensemble fails as the volume grows. Then address a tempting escape route: for two mass-degenerate flavours one may write the measure as $\det M(\mu)\det M(\mu)^\dagger = |\det M(\mu)|^2 \geq 0$, which is manifestly positive and cheap to simulate. What theory has one actually simulated, at what chemical potential for each species, and why does its positivity not constitute a solution?
\end{question}
\begin{solution}
At $\mu = 0$ the fermion matrix obeys $\gamma_5 M \gamma_5 = M^\dagger$, which forces the eigenvalues to come in complex-conjugate pairs and $\det M$ to be real (and positive for degenerate pairs), so it can serve as a probability weight. The finite-density coupling $e^{\pm\mu a}$ breaks this: instead $\gamma_5 M[U,\mu]\gamma_5 = M[U,-\mu^*]^\dagger$, so for real $\mu$ one has $\det M[U,\mu] = (\det M[U,-\mu])^*$ --- relating the determinant to the conjugate of a \emph{different} matrix, not to its own. The determinant is therefore genuinely complex and importance sampling has no positive weight. One can reweight, writing $\langle O\rangle_{\rm full} = \langle O e^{i\varphi}\rangle_{\rm pq}/\langle e^{i\varphi}\rangle_{\rm pq}$, but the denominator is $\langle e^{i\varphi}\rangle_{\rm pq} = Z_{\rm full}/Z_{\rm pq} = e^{-\Omega\,\Delta f}$ with $\Delta f\geq0$, exponentially small in the four-volume. The statistical error of the denominator does not shrink with $\Omega$, so the required statistics grows like $e^{2\Omega\Delta f}$: reweighting fails exponentially fast.

The escape route fails because $|\det M(\mu)|^2$ is not the measure of the theory one wanted. Using $\gamma_5$-Hermiticity, $\det M(\mu)^\dagger = (\det M(\mu))^* = \det M(-\mu)$, so
\begin{equation}
|\det M(\mu)|^2 = \det M(\mu)\,\det M(-\mu).
\end{equation}
That is the measure of a two-flavour theory in which one flavour carries chemical potential $+\mu$ and the other $-\mu$: a theory at finite \emph{isospin} chemical potential $\mu_I = \mu$ with zero baryon chemical potential, not two flavours at a common baryon chemical potential $\mu$. This is exactly the phase-quenched theory. Its positivity is genuine, but its physics is different: the conjugate quark pair can form a light charged-pion-like state with a nonzero net isospin charge, so the theory undergoes pion condensation already at $\mu \approx m_\pi/2$, whereas real QCD at zero temperature stays completely inert until $\mu$ reaches roughly $m_N/3$. The two theories therefore have different free energies --- which is precisely the statement $\Delta f > 0$, i.e.\ the reason the average phase is exponentially small. Positivity has not been bought for free: the entire difficulty has been relocated into the reweighting factor connecting the two theories. This also gives the sharp conceptual statement: the sign problem is the price of the mismatch between the theory one can sample and the theory one wants.
\end{solution}

\begin{question}
\diff{Hard} Three sharp conceptual points about what the sign problem is and is not.
\begin{enumerate}[label=(\roman*), leftmargin=*]
  \item On any single configuration $\det M[U,\mu]$ is complex, yet $Z(\mu)$ is real and moreover even in $\mu$. Identify the two exact properties of the lattice theory responsible for these two statements, and explain why a symmetry that makes $Z$ real does \emph{not} make the integrand real configuration by configuration. Explain why precisely this gap --- cancellation \emph{between} configurations rather than within one --- is what makes the Monte Carlo estimate exponentially expensive.
  \item Without performing the computation, argue conceptually why the naive coupling $\mu\,\bar\psi\gamma_4\psi$ generates a $\mu^2/a^2$ divergence while multiplying the temporal hoppings by $e^{\pm\mu a}$ does not. Base your argument on the fact that $\mu$ enters as the fourth component of a constant \emph{imaginary} abelian gauge field, on how gauge fields couple to lattice fermions, and on what a quark accumulates when it propagates once around the temporal direction. What would go wrong if one instead chose the arithmetically simpler $f(\mu a) = 1 + \mu a$, which also has the correct $\mu=0$ limit and the correct slope at the origin?
  \item Using the chain of inequalities $Z_{\rm full} \le Z_{\rm sq} \le Z_{\rm pq}$ established in the lecture, argue that ``the severity of the sign problem'' is not a property of the physical theory alone but of the positive reference ensemble one chooses to sample. What would a representation with no sign problem look like, and how should one therefore interpret a statement such as ``QCD at finite baryon density has an exponentially hard sign problem''?
\end{enumerate}
\end{question}
\begin{solution}
\textbf{(i)} The two exact properties are:
\begin{itemize}
  \item \emph{Charge conjugation.} The Haar measure is invariant under $U_\mu(n) \to U_\mu^*(n)$ and the gauge action satisfies $S_g[U^*] = S_g[U]$, while the fermion determinant obeys $\det M[U,\mu] = (\det M[U^*,\mu^*])^*$. For real $\mu$ this gives $\det M[U,\mu] = (\det M[U^*,\mu])^*$, so the contributions of the pair $(U, U^*)$ are complex conjugates of each other and $Z(\mu) = Z(\mu)^*$ is real. Equivalently one may replace $\det M \to \mathrm{Re}\det M$ in the integrand without changing $Z$.
  \item \emph{$\gamma_5$-Hermiticity.} $\gamma_5 M[U,\mu]\gamma_5 = M[U,-\mu^*]^\dagger$ gives $\det M[U,\mu] = (\det M[U,-\mu])^*$ on the \emph{same} configuration, so $Z(\mu) = Z(-\mu)^*$; combined with reality this yields $Z(\mu) = Z(-\mu)$, i.e.\ $Z$ is even in $\mu$ (the statement underlying the Taylor expansion in $(\mu/T)^{2n}$).
\end{itemize}
The reason neither makes the integrand real is that both relations map a configuration to a \emph{different} configuration or a different parameter: $U \to U^*$ in the first case, $\mu \to -\mu$ in the second. Reality of $Z$ is a statement about the integral, achieved by pairwise cancellation of imaginary parts between $U$ and its conjugate partner $U^*$; it does not say that the weight on any given $U$ is real, and generically $\mathrm{Im}\det M[U,\mu]\neq0$. Only if a symmetry acted on each configuration \emph{individually} --- as $\gamma_5$-Hermiticity does at $\mu = 0$, mapping $M$ to $M^\dagger$ on the same $U$ --- would positivity hold configuration by configuration.

Why this gap is fatal: Monte Carlo importance sampling can only realize cancellations \emph{statistically}, by sampling configurations and letting their contributions cancel in the average. When the cancellation is intra-configuration (as at $\mu=0$) there is nothing to cancel and the weight is a probability. When it is inter-configuration, the estimator is a difference of two large, nearly equal quantities: the positive and negative (or opposite-phase) contributions each have magnitude of order $Z_{\rm pq}$, while their difference is $Z_{\rm full} = e^{-\Omega\Delta f}Z_{\rm pq}$. The \emph{signal} cancels exponentially but the \emph{noise} --- which adds in quadrature and is set by the typical magnitude of the terms, not their sum --- does not. Hence the relative error scales as $e^{\Omega\Delta f}/\sqrt{N_{\rm cfg}}$ and one needs $N_{\rm cfg}\sim e^{2\Omega\Delta f}$ samples. This is the toy-model picture: $e^{-\phi^2}\cos(\mu\phi)$ has $\mathcal{O}(1)$ positive and negative lobes whose net integral is $\sqrt\pi e^{-\mu^2/4}$; the answer is a tiny residue of a near-perfect cancellation, and each Monte Carlo sample carries $\mathcal{O}(1)$ noise.

\textbf{(ii)} The argument runs through gauge invariance and compactness. On the lattice, gauge fields couple to fermions only through the link variables multiplying the hopping terms --- i.e.\ as \emph{phases} $e^{iaA_\mu}$, never as an additive term $A_\mu\bar\psi\gamma_\mu\psi$ at a site. This is not decoration: it is what makes the coupling of the gauge field to the fermion \emph{finite}, because a link is a bounded (compact) group element and its effect on the propagator is a reweighting of hopping amplitudes rather than the insertion of a new dimensionful operator. Since $\mu$ enters the grand-canonical trace exactly as the fourth component of a constant imaginary abelian gauge potential, $A_4 = -i\mu$, consistency demands that it appear in the same way: on the temporal links, as $f(\mu a) = e^{iaA_4}|_{A_4 = -i\mu} = e^{\mu a}$ forward and $e^{-\mu a}$ backward. Adding $\mu\bar\psi\gamma_4\psi$ as a separate site term instead treats $\mu$ as a \emph{non-compact} additive vector potential; this is an operator insertion that mixes with the cutoff and, since $\mu$ has mass dimension one and the only other scale available is $1/a$, produces additive divergences $\propto \mu^2/a^2$ in the energy density. Equivalently: the naive term is not the minimal coupling of a gauge field to a lattice fermion, and only the minimal (link) coupling is protected.

The complementary way to see that $e^{\mu a}$ is right is the physical statement about winding once around the temporal direction: a quark hopping one step in Euclidean time picks up $e^{\mu a}$, so after $N_\tau$ steps around the periodic time direction it accumulates $e^{\mu a N_\tau} = e^{\mu/T}$ --- exactly the grand-canonical Boltzmann factor $e^{\mu N/T}$ per unit of quark number. The chemical potential must enter through the temporal holonomy, and the exponential is what makes the holonomy come out right.

With $f(\mu a) = 1 + \mu a$: this satisfies $f(0)=1$ and $f'(0)=1$, but it violates the third condition, $f(\mu a) = 1/f(-\mu a)$, since $1/f(-\mu a) = (1-\mu a)^{-1} = 1 + \mu a + (\mu a)^2 + \dots \neq 1 + \mu a$. The consequences are concrete. First, particles and antiparticles are no longer weighted reciprocally, so the exchange $\mu \to -\mu$ no longer corresponds to exchanging forward and backward hopping; the exact relation $Z(\mu) = Z(-\mu)$ and the $\gamma_5$-Hermiticity relation $\det M[U,\mu] = (\det M[U,-\mu])^*$ are spoiled at $\mathcal{O}((\mu a)^2)$. Second and more seriously, winding around the temporal direction gives $(1+\mu a)^{N_\tau} = e^{N_\tau \ln(1+\mu a)} = e^{\mu/T}\,e^{-N_\tau(\mu a)^2/2 + \dots}$, so the accumulated factor is \emph{not} $e^{\mu/T}$; the error per winding is $N_\tau (\mu a)^2 = \mu^2 a/T$, and quantities such as the energy density pick up the very $\mu^2/a^2$-type additive divergence one was trying to avoid. Only a function satisfying all three conditions --- of which $e^{\mu a}$ is the natural representative, being the actual holonomy --- gives a finite continuum limit. (Any $f$ agreeing with $e^{\mu a}$ up to $\mathcal{O}((\mu a)^3)$ and satisfying the reciprocity condition also works, which is why the choice is natural rather than unique.)

\textbf{(iii)} The lecture's inequalities follow from $|\mathrm{Re}\det M| \le |\det M|$ pointwise, giving
\begin{equation}
\langle \mathrm{sign}\rangle_{\rm sq} = \frac{Z_{\rm full}}{Z_{\rm sq}} \;\ge\; \frac{Z_{\rm full}}{Z_{\rm pq}} = \langle e^{i\varphi}\rangle_{\rm pq},
\qquad (\Delta f)_{\rm sq} \le (\Delta f)_{\rm pq}.
\end{equation}
Read this carefully: the numerator $Z_{\rm full}$ is the physical theory and is fixed, but the denominator is the partition function of whichever positive ensemble one decided to sample. Two different but equally legitimate choices --- sign quenching and phase quenching --- give two \emph{different} values of the average reweighting factor and hence two different ``severities''. The severity is therefore not an observable of QCD at finite density; it is a property of the pair (target theory, sampled proxy). The general statement is that for any positive measure $p[U]$ used as a reference, the reweighting cost is governed by $Z_{\rm full}/Z_p$, i.e.\ by the free-energy difference between target and proxy, and nothing in the physics fixes that difference --- it is fixed by the choice of $p$, which includes the choice of \emph{integration variables}.

A representation with no sign problem would be one in which, after an exact change of variables (or an exact resummation of the original degrees of freedom), the weight of every configuration of the \emph{new} variables is real and non-negative, so that $Z_p = Z_{\rm full}$ and $\Delta f = 0$ identically. Such representations are known and used: dual/worldline and flux representations of various lattice models trade the original fields for integer-valued worldline or flux variables in which the finite-density weight is manifestly positive, and simulations there are as easy at $\mu \ne 0$ as at $\mu = 0$. Nothing in the definition of $Z$ obstructs this in principle --- $Z$ is just a number, and a number does not have a sign problem.

Hence the statement ``QCD at finite baryon density has an exponentially hard sign problem'' should be read as shorthand for: ``in the standard representation, in terms of gauge links with the fermions integrated out, and using the standard positive proxies (phase- or sign-quenched), the reweighting factor is exponentially small in the four-volume, and no positive representation of lattice QCD is currently known.'' It is a statement about the state of our formulations, not a theorem about the physics. Two important caveats keep this from being empty optimism: methods like thimbles and complex Langevin are attempts to construct better contours/measures rather than better representations, and general complexity-theoretic results (the sign problem for generic Hamiltonians is NP-hard) make it very unlikely that a universal, always-applicable positive representation exists. The correct attitude is that the sign problem is representation dependent and therefore worth attacking by changing variables, but that no change of variables is guaranteed to exist for a given theory.
\end{solution}

\section{Techniques to circumvent the sign problem}

\begin{question}
\diff{Easy} Distinguish sign quenching from phase quenching. Why is the phase-quenched theory generally not the theory one wants to solve?
\end{question}
\begin{solution}
Phase quenching replaces the complex weight $\det M = |\det M| e^{i\varphi}$ by its modulus $|\det M|$ and corrects with a continuous phase $e^{i\varphi}\in U(1)$; sign quenching uses $|\mathrm{Re}\det M|$ and corrects with a discrete sign $\pm1 \in \mathbb{Z}_2$. Both are ratio reweightings of the same form $\langle O\rangle_{\rm full} = \langle w O\rangle_{\rm proxy}/\langle w\rangle_{\rm proxy}$. In QCD phase quenching corresponds to giving the second flavour the opposite chemical potential, i.e.\ to simulating an isospin rather than a baryon chemical potential. That theory has genuinely different physics --- notably charged pion condensation at $\mu \approx m_\pi/2$ --- and hence a different free energy, so it does not reproduce the finite-baryon-density results one wants: the discarded phase carries the essential physics.
\end{solution}

\begin{question}
\diff{Intermediate} Describe the Taylor expansion about $\mu/T=0$ and the analytic continuation from imaginary chemical potential, stating for each its domain of validity and its main limitation. Both reconstruct a function of $\mu/T$ that contains only \emph{even} powers: state the two exact properties of the lattice action that guarantee this, and explain what would go wrong in the analytic continuation if the Roberge--Weiss periodicity of $Z$ at imaginary $\mu$ were ignored when choosing the fit ansatz.
\end{question}
\begin{solution}
\emph{Taylor expansion:} simulate at $\mu = 0$, where the weight is positive, and expand the pressure in powers of $\mu/T$, $p/T^4 = \sum_n c_{2n}(T)(\mu/T)^{2n}$, computing the coefficients as $\mu=0$ expectation values (generalized quark-number susceptibilities). Valid for small $\mu/T$; limited by the exploding statistical noise of the high-order coefficients (many noisy disconnected contributions; beyond $n\sim4$ the signal is typically lost) and by the finite radius of convergence set by the nearest singularity in the complex $\mu/T$ plane.

\emph{Imaginary chemical potential:} for $\mu = i\mu_I$ one has $\gamma_5 M(i\mu_I)\gamma_5 = M(i\mu_I)^\dagger$, so $\det M$ is real (positive in the cases of practical interest), there is no sign problem, and one simulates directly and then continues analytically to real $\mu$. Valid within the first Roberge--Weiss sector, $\mu_I/T < \pi/3$; limited by the systematic uncertainty of the continuation, i.e.\ by the choice of fit ansatz, and by the RW transition bounding the accessible region.

The evenness in $\mu/T$ follows from the two exact properties established in the lecture: (i) $\gamma_5$-Hermiticity, $\gamma_5 M(U,\mu)\gamma_5 = M(U,-\mu^*)^\dagger$, giving $Z(\mu) = Z(-\mu^*)^*$; and (ii) charge conjugation, $U_\mu(n)\to U_\mu^*(n)$, giving $Z(\mu) = Z(\mu^*)^*$, i.e.\ $Z$ real for real $\mu$. Combining them yields $Z(\mu) = Z(-\mu)$, so $p$ and hence the whole series contains only even powers, and the density is odd. This is also what makes the imaginary-$\mu$ route useful: continuing $\mu \to i\mu_I$ simply alternates the signs of the coefficients, $a_{2n}(\mu/T)^{2n} \to (-1)^n a_{2n}(\mu_I/T)^{2n}$, so a polynomial fit at imaginary $\mu$ determines the same $a_{2n}$ that the Taylor method computes at $\mu = 0$.

If the Roberge--Weiss periodicity $Z(\mu_B + i2\pi T) = Z(\mu_B)$ were ignored, one would be fitting a periodic (and, above $T_{\rm RW}$, non-analytic) function of $\mu_I$ with a polynomial over a range where that is invalid. Two things go wrong. First, the data at imaginary $\mu$ carry no new information beyond one period, so extending the fit range in the hope of better-constrained coefficients is illusory --- and if the range crosses $\mu_I/T = \pi/3$ one crosses the RW transition, where the analytic continuation of the low-$\mu_I$ branch simply does not extend. Fitting through it mixes two different analytic branches and produces coefficients with no relation to the real-$\mu$ physics. Second, even inside the first sector, a fit ansatz inconsistent with the known analytic structure (for example a polynomial where a Fourier ansatz in $\mu_I/T$ is appropriate in the confined phase, where the free energy is dominated by hadronic --- hence $2\pi$-periodic in $\mu_B/T$ --- excitations) biases the extracted $a_{2n}$ in a way that is invisible in the fit quality but fatal for the extrapolation. This ansatz dependence, not statistics, is the dominant systematic of the method.
\end{solution}

\begin{question}
\diff{Hard} What the standard toolbox does and does not achieve.
\begin{enumerate}[label=(\roman*), leftmargin=*]
  \item The Taylor-expansion method and analytic continuation from imaginary $\mu$ both end up determining the same coefficients $a_{2n}$ of the same function. In what sense are they therefore the same computation, and in what sense do they carry genuinely different systematic errors? Explain why neither can reliably locate a critical endpoint, and be explicit about which failure is statistical and which is a matter of principle.
  \item Distinguish the \emph{sign problem} from the \emph{overlap problem}. Using the lecture's toy model $Z(\mu) = \int d\phi\, e^{-\phi^2 + i\mu\phi}$, describe (a) a situation in which the sampled and target distributions overlap essentially perfectly and yet the reweighting factor is still exponentially small, and (b) a situation in which the sampling weight is positive and easy to sample but the overlap is catastrophic. What does this pair of examples tell you about using the average phase as a diagnostic, and about which of the two problems a new method must actually solve?
  \item On a single Lefschetz thimble the oscillating factor $e^{-iS_I}$ is a constant and can be pulled out of the integral --- yet the sign problem is not solved. Enumerate the residual sources of cancellation and say which of them is a property of the \emph{geometry} of the deformed contour and which is a property of the \emph{parametrization} one chooses for it. Finally, explain why the Stokes phenomenon is a genuine obstruction of the method rather than a numerical nuisance that a better integrator would remove.
\end{enumerate}
\end{question}
\begin{solution}
\textbf{(i)} They are the same computation in the following precise sense: both reconstruct the Taylor coefficients $a_{2n}(T)$ of the \emph{same} analytic function $p(\mu/T)/T^4$ about $\mu = 0$, from data taken in a region where the sign problem is absent. The Taylor method obtains them as derivatives at the point $\mu = 0$ (measured as susceptibilities), the imaginary-$\mu$ method as fit parameters from the function's values along the imaginary axis, where by evenness $a_{2n}(\mu/T)^{2n} \to (-1)^n a_{2n}(\mu_I/T)^{2n}$. Since a function is determined by its derivatives at a point or by its values on a segment, the information content is formally the same, and their results for the low-order coefficients agree within errors --- which is exactly why the comparison is used as a cross-check.

The systematics differ genuinely, however, because the two methods reach the coefficients by different routes.
\begin{itemize}
  \item Taylor: the $n$-th coefficient requires a $2n$-th derivative, i.e.\ a high-order cumulant estimated from many noisy disconnected pieces. This is a \emph{statistical} failure: the error on $a_{2n}$ grows rapidly with $n$ and the coefficients drown in noise beyond $n\sim4$. It is unbiased in principle --- with enough statistics one gets the coefficient right.
  \item Imaginary $\mu$: the coefficients come from a \emph{fit} over a finite window $\mu_I/T < \pi/3$ with a truncated ansatz. The statistics are excellent, but the result is biased by the choice of ansatz and by the truncation order, and the bias is not visible in the $\chi^2$. This is a \emph{systematic}, ansatz-dependence failure, and it does not go away with more statistics; it also cannot use data beyond the RW transition, which is not a technical limit but a boundary of analyticity.
\end{itemize}
So the two methods are complementary in a useful way: they share the target but fail differently, and agreement between them is real evidence.

Neither can locate a critical endpoint, and it is worth separating the two reasons. The failure of the Taylor method to get high coefficients is statistical; the failure to reach a critical endpoint is a matter of \emph{principle}. A critical endpoint is a genuine singularity of the free energy at some $\mu_{\rm CEP}$; the Taylor series about $\mu=0$ has a radius of convergence bounded by the distance to the nearest singularity in the complex $\mu/T$ plane, so a truncated series simply cannot represent the function at or beyond that point --- no amount of statistics helps, and low-order truncations extrapolated past the radius produce artefacts that can easily be mistaken for a transition. The same applies to analytic continuation: an analytic function continued from a segment cannot be continued \emph{through} a singularity, and estimates of the radius of convergence (e.g.\ from ratios of coefficients or from the location of RW-related singularities) are themselves limited by exactly the coefficients one cannot determine. Both methods are therefore intrinsically restricted to the small-$\mu/T$ region, which is why they cannot settle the existence or location of the QCD critical endpoint.

\textbf{(ii)} The two problems are logically independent failure modes of reweighting $\langle O\rangle_{\rm full} = \langle wO\rangle_p/\langle w\rangle_p$.
\begin{itemize}
  \item The \emph{sign problem} is the statement that $\langle w\rangle_p$ is a tiny residue of cancellations, so the ratio is $0/0$ up to noise: the signal cancels exponentially while the noise does not.
  \item The \emph{overlap problem} is the statement that the region of configuration space dominating the target measure is not the region sampled by $p$: the reweighting is then controlled by rare configurations in the tail of $p$, and the estimate is dominated by a handful of samples with enormous weights (large variance, badly biased in practice) even if no cancellation were involved at all.
\end{itemize}
In the toy model, $Z(\mu) = \sqrt\pi e^{-\mu^2/4}$ and the integrand is $e^{-\phi^2}\cos(\mu\phi)$ (real part).
\begin{itemize}
  \item[(a)] \emph{Good overlap, exponentially small reweighting factor.} Take the sign-quenched reference $p \propto |e^{-\phi^2}\cos(\mu\phi)|$. By construction it has exactly the same support and the same envelope as the full integrand --- the two curves differ only by the sign of alternating lobes, as in the lecture's figure --- so the overlap is essentially perfect and no rare configurations are involved. Nevertheless the average sign is $\langle \mathrm{sign}\rangle_{\rm sq} = Z(\mu)/Z_{\rm sq}(\mu)$, and since $Z_{\rm sq} = \int d\phi\, e^{-\phi^2}|\cos\mu\phi| = \mathcal{O}(1)$ for large $\mu$ while $Z = \sqrt\pi e^{-\mu^2/4}$, the reweighting factor is still exponentially small. The cancellation between lobes is intrinsic to the target; sampling it perfectly does not help.
  \item[(b)] \emph{Positive, easily sampled weight, catastrophic overlap.} Take the phase-quenched reference $p \propto e^{-\phi^2}$, a plain Gaussian --- trivially positive and trivially sampled. But the full integrand is concentrated, in the sense that matters, around the shifted saddle $\phi = i\mu/2$ (completing the square gives $Z = \sqrt\pi e^{-\mu^2/4}$ with the stationary point off the real axis), and the misalignment between the Gaussian and the oscillating integrand grows with $\mu$: this is exactly the lecture's figure showing the phase-quenched weight and the full integrand drifting apart. Here \emph{both} problems are present, and $\langle e^{i\varphi}\rangle_{\rm pq} = e^{-\mu^2/4}$ is smaller than the sign-quenched factor --- consistent with $Z_{\rm full}\le Z_{\rm sq}\le Z_{\rm pq}$.
\end{itemize}
The lessons: first, the average phase is a diagnostic of the \emph{pair} (target, proxy) and not of the target --- it measures a free-energy difference, so a small value can mean either that the target has intrinsic cancellations or merely that a poor proxy was chosen, and comparing $\langle\mathrm{sign}\rangle_{\rm sq}$ with $\langle e^{i\varphi}\rangle_{\rm pq}$ separates the two. Second, and more important for method development: improving overlap alone is not enough. Example (a) shows a method with perfect overlap that still fails; so a successful method must eliminate the \emph{cancellation}, i.e.\ produce a genuinely positive measure (a positive representation, or a contour on which the phase is constant), not merely a better-aligned positive proxy. This is precisely the motivation for contour deformation and dual representations rather than smarter reweighting.

\textbf{(iii)} The oscillation is removed only \emph{within} a single thimble; the residual sources of cancellation are:
\begin{enumerate}[label=(\alph*), leftmargin=*]
  \item \textbf{Inter-thimble interference.} The original real cycle decomposes as $\mathbb{R}^N = \sum_\sigma m_\sigma \mathcal{J}_\sigma$ with integer intersection numbers $m_\sigma$, so
  $Z = \sum_\sigma m_\sigma e^{-iS_{I,\sigma}}\int_{\mathcal{J}_\sigma}\mathcal{D}\phi_c\,e^{-S_R}$. Each thimble carries its own constant phase $e^{-iS_{I,\sigma}}$ and its own sign from $m_\sigma$, so when several thimbles contribute the sum can cancel --- a global sign problem, though usually much milder, since the number of contributing saddles is small compared with the dimension of field space. This is a property of the \emph{geometry}: which saddles exist, which thimbles the real cycle intersects, and with what multiplicities. It is fixed by the action and the parameters, and cannot be improved away by a better choice of coordinates. (Determining the $m_\sigma$ is itself one of the main unsolved practical problems.)
  \item \textbf{The residual phase from the Jacobian.} A thimble is a real $N$-dimensional manifold embedded in $\mathbb{C}^N$, and to integrate on it one must parametrize it --- in practice by flowing from a tangent-space basis at the saddle. The induced measure carries the Jacobian of that map, $\det J$, which is in general complex and varies over the thimble, so $e^{-S_R}|\det J|$ is the sampling weight and $\arg\det J$ must be reweighted. This is the \emph{residual sign problem}, and it is a property of the \emph{parametrization}: the geometric content (the orientation of the thimble) is fixed, but how much phase fluctuation one has to fight depends on the coordinates chosen, and it can be reduced by better parametrizations --- which is exactly what the generalized-thimble and learned-contour methods exploit. Empirically these fluctuations are mild compared with the original oscillations, but they are not zero, and computing $\det J$ at every step is the dominant numerical cost.
  \item One should add the practical cost of integrating the flow itself, and the fact that thimbles can terminate at zeros of the fermion determinant, where $S_{\rm eff} = S_g - \ln\det M$ has logarithmic branch points --- so the flow must be defined with the holomorphic weight and the contour can be obstructed.
\end{enumerate}
The Stokes phenomenon is a genuine obstruction because it is a statement about the \emph{decomposition}, not about numerics. As parameters (e.g.\ $\mu$, or the couplings) are varied, one can reach a configuration in which the descent manifold of one saddle flows directly into another saddle; on such a Stokes ray, $\mathrm{Im}\,S$ coincides for the two saddles, the thimbles are not transverse to the anti-thimbles, and the intersection numbers $m_\sigma$ are \emph{ill-defined} --- they jump discontinuously as one crosses. The decomposition of $\mathbb{R}^N$ into thimbles therefore changes discontinuously, so a set of contributing thimbles determined on one side of the ray is simply the \emph{wrong} answer on the other side (this is why the exercise instructs one to use different decompositions on the two sides of $\sigma_R = 0$). No integrator can fix this: the underlying integral $Z$ is perfectly continuous, but the representation used to compute it is not, and one must know the topology --- which saddles contribute and with what multiplicity --- before integrating. In a realistic lattice theory, with a high-dimensional complexified field space and a large number of saddles whose intersection numbers are not computable in practice, one cannot in general certify that one is on the correct side of every Stokes ray. That, more than the residual phase, is why the thimble programme has not displaced reweighting and analytic continuation as the workhorse for finite-density lattice QCD.
\end{solution}

\end{document}
